\begin{tikzpicture}[
  x=1mm,y=1mm,
  >=Stealth,
  font=\figurefont\fontsize{8}{10}\selectfont
]

\tikzset{
  vneg/.style={circle,draw=FigureBlue,fill=FigureBlueFill,minimum size=4.8mm,inner sep=0pt,line width=.6pt,text=FigureInk},
  vpos/.style={circle,draw=FigureGreen,fill=FigureGreenFill,minimum size=4.8mm,inner sep=0pt,line width=.6pt,text=FigureInk},
  edge/.style={draw=FigureStroke,line width=.65pt},
  bridge/.style={draw=FigureRed,line width=.95pt},
  cutline/.style={dashed,draw=FigureRed,line width=.65pt},
  note/.style={font=\FigureNoteFont\fontsize{7.5}{9}\selectfont,text=FigureInk},
  muted/.style={font=\FigureNoteFont\fontsize{7.2}{8.8}\selectfont,text=FigureMuted}
}

% Left panel: Barbell graph topology (two triangles connected by bridge 3-4)
\begin{scope}[xshift=0mm,yshift=4mm]
  \node[vneg] (v1) at (0,12) {$1$};
  \node[vneg] (v2) at (0,0) {$2$};
  \node[vneg] (v3) at (9,6) {$3$};
  \node[vpos] (v4) at (21,6) {$4$};
  \node[vpos] (v5) at (30,12) {$5$};
  \node[vpos] (v6) at (30,0) {$6$};

  \draw[edge] (v1)--(v2)--(v3)--(v1);
  \draw[edge] (v4)--(v5)--(v6)--(v4);
  \draw[bridge] (v3)--(v4);
  \draw[cutline] (15,-3) -- (15,15);

  \node[muted] at (4,16.5) {$(\bm v_2)_i<0$};
  \node[muted] at (26,16.5) {$(\bm v_2)_i>0$};
  \node[note,text=FigureRed] at (15,-5.5) {瓶颈桥边 $3\!-\!4$};
  \node[note] at (15,-11) {(a) 杠铃图的正负节域};
\end{scope}

% Right panel: 1D spring embedding along Fiedler vector v_2
\begin{scope}[xshift=42mm,yshift=4mm]
  % Axis
  \draw[->,draw=FigureStroke,line width=.6pt] (0,6) -- (54,6) node[right,note] {$(\bm v_2)_i$};
  \draw[cutline] (27,-2) -- (27,15);
  \node[note,text=FigureRed] at (27,16.8) {$0$（节线切点）};

  \node[vneg] (e1) at (7,10.5) {$1$};
  \node[vneg] (e2) at (7,1.5) {$2$};
  \node[vneg] (e3) at (17,6) {$3$};
  \node[vpos] (e4) at (37,6) {$4$};
  \node[vpos] (e5) at (47,10.5) {$5$};
  \node[vpos] (e6) at (47,1.5) {$6$};

  \draw[edge] (e1)--(e2);
  \draw[edge] (e1)--(e3);
  \draw[edge] (e2)--(e3);
  \draw[bridge] (e3)--(e4);
  \draw[edge] (e4)--(e5);
  \draw[edge] (e4)--(e6);
  \draw[edge] (e5)--(e6);

  \node[muted] at (11,-3.5) {左团紧聚};
  \node[muted] at (43,-3.5) {右团紧聚};
  \node[note] at (27,-11) {(b) Fiedler向量一维弹簧嵌入};
\end{scope}
\end{tikzpicture}
